\subsection{Symbolic Program Construction and Correspondence}
\label{app:symbolic-compilation-proof}
We prove Theorem~\ref{thm:symbolic-compilation} for admitted plans in
Section~\ref{sec:formal-scope}. Admission checks the operation templates,
read-once components, matching input sets, captured source fields, and
independent unrestricted RHS fields. The proof starts after frontend
preparation; it does not decide whether arbitrary operation declarations
have equivalent meanings.

\paragraph{Meaning of the program.}
The program keeps axis membership separate from axis order. Sets are used
only in set operations, membership tests, and as the membership argument
of an ordered filter. Their enumeration order therefore cannot change an
attribute. Filters preserve their input order; concatenation joins ordered
lists. Position lookup requires a distinct operand axis list and distinct
queried axes, all present in that operand. Each expression refers only to
earlier expressions, whose values can be reused. All component checks are
evaluated before any attribute builder; the checks use no partial position
lookup.

\begin{proof}[Proof of Theorem~\ref{thm:symbolic-compilation}]
We establish three invariants by following the compiler's passes.

\noindent\textbf{1. Source expressions represent the same axes and orders.}
Replay follows Section~\ref{sec:axis-model}. Inputs introduce distinct
ports. A permutation reorders them; a sum removes selected output axes
without changing their identities. A contraction joins only the declared
batch and contraction partners, checks paired extents, and records the
batch/left-free/right-free output order. Thus induction over source
operations gives the same logical axes and ordered interfaces as the
mathematical model. Equal extents alone never identify axes. Additive
components keep separate axis identities; source addition checks their
ordered shapes.

Mappings and partitions preserve this information. Gather composition
selects the same entry $Q[q[h[j]]]$ on both sides of
Equation~\eqref{eq:gather-composition}. In paired matching,
$\gath(d_\ell,u_G)=\gath(q_G,v_G)$ selects the same axes, and applying
$u_G$ to $d_r$ keeps their contraction partners. Selected positions and
their ordered complement cover the group without overlap. Their orders
are recorded explicitly; concatenating them need not restore the original
order. These facts hold for every finite group length.

\noindent\textbf{2. Generated checks equal the mathematical checks.}
The inputs inside each RHS subtree are known at compile time, so
Equation~\eqref{eq:emitted-mandatory} computes exactly $M_v$. Induction over
the forward pass gives the same retained sets $S_v$: inputs keep their
axes, permutations keep the child's set, sums intersect it with $M_v$,
and contractions use Equation~\eqref{eq:canonical-roles}.

The reverse pass starts with $F$ and computes the same $Req_v$ as
Section~\ref{sec:ordering-feasibility}. Distinctness and inclusion check
that requirements contain distinct available axes. At inputs, filtering
the physical axis order by the required set and comparing with $Req_v$
implements the subsequence check. At contractions, the check compares
$Req_v$ with its batch/left-free/right-free parts in that order. The root
check compares $S_{\mathrm{root}}$ with $\operatorname{set}(F)$.
These total tests use only source information and the retained sets,
never attribute builders. Their conjunction is
$\operatorname{check}(p)$, proving Equation~\eqref{eq:emitted-check} even
when construction is infeasible.

\noindent\textbf{3. Successful checks give the same attributes.}
Assume the checks pass. Induct over the RHS operations using
Section~\ref{sec:formal-construction}'s construction. Inputs use their
source orders. A permutation places $Req_v$ first and appends the other
axes in child order. A sum selects the removed axes and preserves the
survivors' order. A contraction extends the required batch order and
chooses remaining batch and contraction axes in left-operand order,
using the same contraction order in both operands. This is the tie rule
$\tau$.

The role sets and checks ensure that every queried axis has a unique
position in its operand; source legality supplies paired extent equality.
Thus the position lookups are defined and produce the same attributes as
$\operatorname{build}_{\tau}(p)$. Each output contains exactly $S_v$ and
satisfies $Req_v$; the root therefore has order $F$. Independent additive
components combine as in Appendix~\ref{sec:addition-proof}. This proves
Equation~\eqref{eq:emitted-build} without reproving Section~\ref{sec:formal-correctness}'s
correctness of the mathematical construction.

Finally, each pass appends finitely many expressions for a finite plan,
including its finite component expansion. Emission needs neither concrete
metadata, the written guard, nor successful simplification. Compilation
therefore terminates with a finite program, subject to the theorem's
resource and representation assumptions.
\end{proof}

\paragraph{Checked symbolic lowering.}
Lowering rewrites this same program into source-attribute expressions.
The mapping identities above and the meaning of each set and list
operation preserve its feasibility checks and positional results.
Membership sets may be reordered because their enumeration order is unused.
Lowering may also reorder a contraction group in both operands by the same
permutation. This changes the integer lists but only reindexes the finite
sum, preserving paired reads, extents, values, and output order.
The independent validator checks the completed candidate and original
metadata expressions under the written guard
(Theorem~\ref{thm:validator-soundness}); unproved construction goals are
never assumptions. Failed lowering gives no expression-based acceptance
and leaves the original program intact. Neither lowering nor guard proof
search is claimed complete. Section~\ref{sec:implementation-scope} describes
the implementation and its trusted components.
